\section{Rappels de théorie des nombres}

 Le stage a débuté par une mise à niveau en théorie des nombres, dont on énonce les résultats principaux utiles par la suite dans cette section. On se réfère à \cite{Cox}[Ch 5, 7].

\subsection{Ordres dans un corps de nombre}

\begin{Def}
	
	Un corps $\K$ est un corps de nombre si c'est une extension finie de $\Q$. Un ordre de $\K$ est un sous anneaux unitaire de $\K$, finement engendré comme $\Z$-module et contenant une $\Q$-base de $\K$.
\end{Def}

\begin{Ex}
	On note $\OR_{\K}$ l'anneau des entiers algébriques d'un corps de nombre $\K$. C'est un ordre de $\K$.
\end{Ex}

\begin{Prop}\cite{Cox}[Prop 5.3, 7.2]
	
	Soit $\OR$ un ordre de $\K$. Alors $\OR$ est un $\Z$-module libre de rang $\left[ \K : \Q\right] $, inclus dans $\OR_{\K}$.
	De plus $\K = Frac(\OR)$
\end{Prop}

$\OR_{\K}$ est donc un ordre maximal. Dans la suite, le cas où $\K$ est une extension quadratique imaginaire de $\Q$ joue un rôle clef. On commence par décrire la structure des ordres dans ce cas:

\begin{Def}
	Un corps $\K$ est un corps quadratique imaginaire si c'est une extension quadratique imaginaire de $\Q$, c'est à dire $\K = \Q(\sqrt{D})$ avec $D$ entier négatif sans facteur carré. 
\end{Def}

Lorsque $\K = \Q(\sqrt{D})$, on peut réécrire $\K = \Q(\sqrt{\Delta})$ où $\Delta = 4D$ est le discriminant de $X^{2} - D$. 

\begin{Def}
	Le discriminant d'une extension quadratique imaginaire $\K = \Q(\sqrt{D})$, noté $d_{\K}$, est défini par :
	\begin{enumerate}
		\item $d_{\K} = D$ si $D = 1 [4]$,
		\item $d_{\K} = 4D$ sinon.
	\end{enumerate}
	En particulier $d_{\K} = 0 ,\;1[4]$. On pose de plus :
	\begin{equation*}
		w_{\K} = \frac{d_{\K} + \sqrt{d_{\K}}}{2}
	\end{equation*}
\end{Def}


\begin{Prop}\cite{Cox}[Prop 5.14, 7.2] 
	
	Soit $\K$ un corps quadratique imaginaire. Alors $\OR_{\K} = \Z\left[ w_{\K}\right] $. De plus si $\K = \Q(\sqrt{D})$ avec $D$ sans facteur carré, alors
	\begin{enumerate}
		\item Si $D \neq 1 [4]$, $\OR_{\K} = \Z\left[ \sqrt{D}\right] $
		\item Sinon, $\OR_{\K} = \Z\left[ \frac{1 + \sqrt{D}}{2}\right] $
	\end{enumerate} 
	
\end{Prop}

On peut maintenant décrire la structure des ordres inclus dans $\OR_{\K}$ :

\begin{Prop}\cite{Cox}[Prop 7.2] Soit $\OR$ un ordre de $\K = \Q(\sqrt{D})$ un corps quadratique imaginaire. $\OR\subset\OR_{\K}$ et on note $f = \left[ \OR_{\K} : \OR\right]$. Alors $\OR = \Z\left[ fw_{\K}\right] $.  
\end{Prop}

\begin{Def}
	Lorsque $\K$ est un corps quadratique imaginaire, si $\OR$ est un ordre de $\K$, $f = \left[ \OR_{\K} : \OR\right]$ est le conducteur de $\OR$, et $\Delta = f^2d_{\K}$ son discriminant.
\end{Def}

La notion de conducteur est multiplicative :

\begin{Prop}
	Soit $\OR$, $\OR'$ deux ordres d'un $\K$ quadratique imaginaire. Soit $f$ et $f'$ leurs conducteurs. Si $\OR'\subset\OR$ alors $f$ divise $f'$
\end{Prop}

\begin{proof}
	D'après la proposition 1.3, $f'w_{\K}\in\OR'\subset\OR$ donc $f'w_{\K}\in Z[fw_{\K}]$. Donc $f'w_{\K} = a + bfw_{k}$ avec $a$, $b\in\Z$. Mais $f'w_{\K}\in\OR_{\K}$ admet une unique écriture dans la $\Z$-base $(1,w_{K})$ de $\OR_{\K}$. Donc $a = 0$ et $f' = bf$. Donc $f$ divise $f'$.
\end{proof}

On peut reformuler ces résultats en terme de réseau de $\C$ :

\begin{Def}
	Un réseau $\Lambda$ de $\C$ est un ensemble de la forme $\left\lbrace nz_{1} + mz_{2} , n,m\in\Z\right\rbrace$ où $z_{1},\; z_{2}\in\C$ sont linéairement indépendant sur $\R$. On note alors $\Lambda = \left[ z_{1}, z_{2}\right] $
\end{Def}

\begin{Rem}
	Les ordres quadratiques imaginaires sont donc des réseaux de $\C$, qui contiennent $1$ : $\OR_{\K} = \left[1 , w_{\K}\right]$, et $\OR = \left[ 1, fw_{\K}\right]$
\end{Rem}

Les ordres d'un corps de nombre $\K$ ne sont pas des anneaux factoriels en général. On va cependant obtenir des résultats de factorisation au niveau des idéaux.

On énonce d'abord les premières propriétés sur les idéaux des ordres de $\K$.

\begin{Prop} Soit $\LR\subset\OR$ un idéal d'un ordre de $\K$.
	\begin{enumerate}
		\item Il existe $a\in\Z$ appartenant à $\LR$
		\item $\LR$ est un $\Z$-module libre de rang $\left[ \K : \Q\right] $
		\item Le quotient $\OR/\LR$ est fini.
		\item Si $\LR$ est premier, alors $\LR$ est maximal. 
	\end{enumerate}
\end{Prop}

\begin{proof}
	On se réfère à \cite{Cox}[Prop 5.4,Ex 5.1]. L'énoncé ne porte que sur $\OR_{\K}$ mais les arguments utilisés sont aussi valable pour un ordre $\OR$ quelconque : voir \cite{Cox}[Ex 7.4]. 
\end{proof}

\begin{Def}
	La norme d'un idéal $\LR\subset\OR$ d'un ordre d'un corps de nombre $\K$ est $N(\LR) = |\OR/\LR|$.
\end{Def}

Dans le cas des corps quadratiques imaginaires, on en déduit la structure des idéaux :

\begin{Coro}
	Soit $\OR$ un ordre de $\K = \Q(\sqrt{D})$ un corps quadratique imaginaire. Tout idéal $\LR$ de $\OR$ est un réseaux de $\C$ de la forme $\left[ a, b + cfw_{\K}\right]$ avec $a = N(\LR)$, $b$, $c\in\Z$, $c\neq 0$ et $0 \leq b< a$. 	
\end{Coro}

\begin{proof}
	D'après la proposition 1.5, $\LR\cap\Z$ est non vide. $\LR\cap\Z$ est un idéal de $\Z$, donc principal engendré par un entier $a > 0$. Alors l'anneau $\OR/\LR$ est de caractéristique $a$, donc $a = N(\LR)$. Soit $\alpha\in\LR$ linéairement indépendant de $a$. Comme $\LR$ est un sous ensemble discret de $\C$, on peut choisir $\alpha$ de norme minimale. 
	
	On montre par l'absurde que $\LR = \left[ a,\alpha\right]$ : soit $z\in\LR$ qui n'est pas dans le réseau $\left[ a,\alpha\right]$. On peut choisir $z$ de norme minimal. Si $z$ est entier, $z\in \LR\cap\Z = a\Z$ ce qui est exclus. Donc $z$ est linéairement indépendant de $a$. Or $|z| \leq |\alpha| + |z - \alpha|$. Si $z\neq\alpha$, $|z| < |\alpha|$, ce qui contredit la minimalité de $|\alpha|$. 
	
	Donc $\LR = \left[ a, \alpha\right]$. Si l'on note $\alpha = b + cfw_{\K}$, $\alpha$ n'étant pas entier on a $c \neq 0$. De plus $\LR = \left[ a, ka + \alpha\right]$ pour tout $k\in\Z$, donc on peut choisir une base telle que $0 \leq b< a$. 
\end{proof}

Pour tout corps $\K$, on montre que $\OR_{\K}$ est un anneaux de Dedekind, un type d'anneaux qui admet une unique factorisation au niveau des idéaux par des idéaux premiers.

\begin{The}\cite{Cox}[Prop 5.5,Cor 5.6]
	
	$\OR_{\K}$ est un anneaux de Dedekind, c'est à dire :
	\begin{enumerate}
		\item $\OR_{\K}$ est intégralement clos
		\item $\OR_{\K}$ est noethérien
		\item Ses idéaux premier sont maximaux 
	\end{enumerate}
	En particulier tout idéal $\LR$ admet une unique factorisation par des idéaux premier (à permutation prêt des facteurs)
	\begin{equation*}
		\LR = \PR_{1}^{e_{1}}\PR_{2}^{e_{2}}\ldots\PR_{g}^{e_{g}}
	\end{equation*}
	avec $e_{1}$, ... $e_{g}\in\N$ et $g\in\N$.
\end{The}

Un ordre $\OR\subset\K$ quelconque n'est pas un anneaux de Dedekind, mais "presque" : seule la clôture intégrale pose problème.

\begin{Prop}\cite{Cox}[5.A, Ex 7.4]
	
	Un ordre $\OR$ d'un corps de nombre $\K$ est noethérien, et ses idéaux premiers sont maximaux.
\end{Prop}

\begin{proof}\cite{Cox}[5.A, Ex 7.4]
	
	Les arguments sont les mêmes que pour $\OR_{\K}$, et découlent de la proposition 1.5.
\end{proof}

Dans la suite, nous allons préciser l'étude de $\OR_{\K}$, tenter de généraliser les résultats obtenues à un ordre quelconque $\OR$, puis déduire un résultat de factorisation d'idéaux d'ordre de corps quadratiques imaginaires. 







\subsection{Idéaux fractionnaires et Groupe de classe}

En plus d'être un anneaux de Dedekind, $\OR_{\K}$ est un anneaux "presque" principal. Pour définir proprement cela, on veut introduire une relation d'équivalence sur les idéaux: $\LR\simeq\LR'\iff\LR = \alpha\LR',\alpha\in\OR_{\K}$. Cette définition est mauvaise car la relation n'est pas à priori symétrique : $\alpha$ n'est pas inversible dans $\OR_{\K}$ en général. Il faut donc considérer $\alpha\in\K$.

On défini donc un notion plus large d'idéaux, les idéaux fractionnaire, pour lesquels cette relation d'équivalence est bien définie :

\begin{Def}
	Un idéal fractionnaire de $\OR_{\K}$ est un $\OR_{\K}$-module de $\K$ de type fini, non nul. On note $I(\OR_{\K})$ l'ensemble des idéaux fractionnaires. Les idéaux de $\OR_{\K}$ seront appelé "idéaux entiers" pour les distinguer des idéaux fractionnaires. 
\end{Def}

\begin{Prop}\cite{Cox}[5.A, Ex 5.2]
	
	Les idéaux fractionnaires sont les ensembles de la forme $\alpha\LR$, où $\alpha\in\K$ et $\LR\subset\OR_{\K}$ est un idéal entier de $\OR_{\K}$.
\end{Prop}

\begin{Rem}
	La définition d'idéaux fractionnaire et la proposition 1.9 précédente se généralise sans soucis aux ordres quelconques. 
\end{Rem}


Le nom d'idéal fractionnaire vient de l'existence d'une structure de groupe multiplicatif sur $I(\OR_{\K})$ : on peut inverser les idéaux "entier". 

\begin{Prop}\cite{Cox}[Prop 5.7]
	
	Les idéaux fractionnaire $\IR$ de $\OR_{\K}$ sont inversible :
	Il existe $\IR'$ tel que $(\IR)(\IR') = \OR_{\K}$.
\end{Prop}

Les idéaux fractionnaire sont suffisamment proche des idéaux entiers pour que l'on obtienne un résultat de factorisation :

\begin{Prop}\cite{Cox}[Prop 5.7]
	
	Soit $\IR$ un idéal fractionnaire de $\OR_{\K}$. Il existe un entier $d$ tel que $\LR = d\IR$ soit un idéal entier. 
\end{Prop}

\begin{Def}
	Si $\IR$ un idéal fractionnaire de $\OR_{\K}$, l'entier $d$ minimal tel que $\LR = d\IR$ soit entier est appelé le dénominateur de $\IR$
\end{Def}

\begin{Prop}\cite{Cox}[Prop 5.7]
	
	Tout idéal fractionnaire $\IR$ admet une unique factorisation par des idéaux entiers premiers ou leurs inverses (à permutation prêt des facteurs):
	\begin{equation*}
		\IR = \PR_{1}^{e_{1}}\PR_{2}^{e_{2}}\ldots\PR_{g}^{e_{g}}
	\end{equation*}
	avec $e_{1}$, ... $e_{g}\in\Z$ et $g\in\N$.
\end{Prop}

On peut étudier les idéaux principaux parmi les idéaux fractionnaires.

\begin{Def}
	Pour tout corps de nombre $\K$, $I(\OR_{\K})$ contient l'ensemble $P(\OR_{\K})$ des idéaux fractionnaires principaux, i.e de la forme $\alpha\OR_{\K}$ avec $\alpha\in\K$. 
\end{Def} 

\begin{Rem}
	L'inverse de $\alpha\OR_{\K}$ est $\alpha^{-1}\OR_{\K}$ : $P(\OR_{\K}))$ est un sous-groupe de $I(\OR_{\K})$. Le quotient $I(\OR_{\K}/P(\OR_{\K}))$ pour la loi multiplicative correspond aux classes d'équivalences de la relation $\IR\simeq\IR'\iff\IR = \alpha\IR',\alpha\in\\K$.
\end{Rem}

\begin{Def}
	Pour tout corps de nombre $\K$, le groupe de classe de $\OR_{\K}$ est le quotient $Cl(\OR_{\K}) =  I(\OR_{\K})/P(\OR_{\K})$
\end{Def}

Le groupe de classe peut être défini pour tout anneaux de Dedekind, en se plaçant dans le corps de fraction. Une propriété remarquable de $\OR_{\K}$ comparé aux anneaux de Dedekind quelconque est que $Cl(\OR_{\K})$ est fini :

\begin{Prop}\cite{Cox}[5.A]
	
	$Cl(\OR_{\K})$ est fini, on note $h(\OR_{\K})$ son cardinal. 
\end{Prop} 

On parlera parfois du groupe de classe de $\K$ pour désigner $Cl(\OR_{\K})$.







\subsection{Groupe de classe d'un ordre quelconque}

On va maintenant définir un groupe de classe pour un ordre $\OR$ quelconque d'un corps de nombre $\K$. on en déduira un résultat de factorisation pour certains idéaux entiers de $\OR$.

Certains idéaux entiers d'un ordre $\OR$ ne seront pas inversible dans $\OR$. Si $\LR\subset\OR$ est un idéal, $\LR$ pourrait aussi être un idéal d'un anneaux contenant $\OR$. 

On a toujours $\OR\subset\left\lbrace \alpha\in\K, \alpha\LR\subset\LR\right\rbrace$. Cependant l'inclusion peut être stricte. On ne pourra définir une notion d'inverse pour $\LR$ qu'en le voyant comme idéal du plus gros anneaux possible. On défini donc la notion d'idéal propre :

\begin{Def}
	$\LR\subset\OR$ est un idéal propre de l'ordre $\OR$ si c'est un idéal entier tel que
	\begin{equation*}
		\OR = \left\lbrace \alpha\in\K, \alpha\LR\subset\LR\right\rbrace
	\end{equation*} 
	$\IR$ est un idéal fractionnaire propre de $\OR$ s'il vérifie $\OR = \left\lbrace \alpha\in\K, \alpha\IR\subset\IR\right\rbrace$. On note $I(\OR)$ l'ensemble des idéaux fractionnaires propres.
\end{Def}

\begin{Prop}\cite{Cox}[7.4]
	
	Un idéal fractionnaire $\IR$ d'un ordre $\OR$ est inversible si et seulement s'il est propre.
\end{Prop}

La restriction aux idéaux propres ne permet pas encore d'obtenir un résultat de factorisation, mais elle permet déjà de définir un groupe de classe.

\begin{Prop}
	Les idéaux fractionnaires principaux d'un ordre $\OR$ sont propres, et forment un sous groupe de $I(\OR)$
\end{Prop}

\begin{proof}
	Les idéaux fractionnaires principaux sont inversibles, donc propre d'après la proposition 1.14. L'ensemble des idéaux fractionnaires principaux est un groupe, et on vient de montrer qu'il est inclus dans $I(\OR)$.
\end{proof}

\begin{Def}
	Soit $\OR$ un ordre d'un corps de nombre. On note $P(\OR)$ le sous groupe de $I(\OR)$ des idéaux fractionnaires principaux. Le groupe de classe de $\OR$ est le quotient $Cl(\OR) =  I(\OR)/P(\OR)$
\end{Def}

\begin{Prop}\cite{Cox}[7.A]
	
	$Cl(\OR)$ est fini, on note $h(\OR)$ son cardinal. 
\end{Prop} 

\begin{Rem}
	On sera amener à chercher des familles génératrices et des factorisations dans le groupe de classes d'ordre $\OR$ de corps $\K$ quadratique imaginaire. Cette étude sera faite plus loin dans le contexte d'un algorithme de calcul d'isogénies 
\end{Rem}







\subsection{Factorisation d'idéaux premier avec le conducteur}


On se concentre désormais sur les ordres $\OR$ de corps quadratiques imaginaires $\K$. On va obtenir un résultat de factorisation de certains idéaux propres ayant des propriétés sur leur norme. Pour cela nous revenons d'abord sur les propriétés de la norme d'un idéal entier.

\begin{Def}
	Soit $\K$ un corps quadratique imaginaire. Il admet un unique automorphisme de corps non trivial. Pour tout $\alpha\in\K$, on note $\bar{\alpha}$ son conjugué dans $\K$. Pour tout idéal $\LR$ d'un ordre $\OR$ de $\K$, on note $\bar{\LR} = \left\lbrace \bar{\alpha}, \alpha\in\LR\right\rbrace $
\end{Def}

\begin{Prop}\cite{Cox}[Prop 7.14]
	
	Soit $\LR$ un idéal entier propre d'un ordre $\OR$ d'un corps quadratique imaginaire $\K$. 
	\begin{enumerate}
		\item $\forall\alpha\in\OR,\; N(\alpha\OR) = N_{\Q}^{\K}(\alpha)$
		\item $N(\LR\LR') = N(\LR)N(\LR')$ pour tout idéal entier propre $\LR'$.
		\item $\LR\bar{\LR} = N(\LR)\OR$, donc $\bar{\LR}$ est l'inverse de $\LR$ dans $Cl(\OR)$
	\end{enumerate}
\end{Prop}

Si $\LR\subset\OR$ est un idéal, le produit $\LR_{\K} = \LR\OR_{\K}$ fourni un idéal de $\OR_{\K}$. Inversement si $\LR_{\K}\subset\OR_{\K}$ est un idéal, l'intersection $\LR = \LR_{\K}\cap\OR$ est un idéal de $\OR$

On va montrer que ces deux applications sont réciproque l'une de l'autre et induisent un isomorphisme entre groupe de classe. 

Mais ce ne sera pas un isomorphisme entre les groupes $I(\OR)$ et $I(\OR_{\K})$. Il faut se restreindre à des sous groupes en imposant une condition sur la norme des idéaux entiers de $\OR$ et de $\OR_{\K}$.

\begin{Def}\cite{Cox}[Lem 7.18]
	Soit $\LR\subset\OR$ un idéal entier, et $f$ le conducteur de $\OR$. $\LR$ est dit premier avec le conducteur si l'une des deux conditions équivalente suivante est vérifiée :
	\begin{enumerate}
		\item $N(\LR)$ est premier avec $f$
		\item $\LR + f\OR = \OR$
	\end{enumerate}
\end{Def}

\begin{Rem}
	La seconde formulation de la définition d'un idéal premier avec le conducteur est analogue à une "relation de Bézout" dans $\OR$ entre le conducteur $f$ et un élément de $\LR$.
\end{Rem}

\begin{Prop}\cite{Cox}[Lem 7.18]
	
	Si $\LR\subset\OR$ est un idéal premier avec le conducteur, alors il est propre. L'ensemble des idéaux premier avec le conducteur est stable par multiplication et forme un sous groupe de $I(\OR)$.  
\end{Prop}

\begin{Def}
	On note $I(\OR,f)$ le sous groupe de $I(\OR)$ engendré par les idéaux premier avec le conducteur, et $P(\OR,f)$ le sous groupe de $I(\OR,f)$ engendré par les idéaux principaux premier avec le conducteur.
\end{Def}

Cette restriction permet tout de même de décrire complètement le groupe de classe. 

\begin{Prop}\cite{Cox}[Lem 7.19]
	
	Si $\OR$ est un ordre d'un corps quadratique imaginaire, de conducteur $f$, alors
	\begin{equation*}
		Cl(\OR) =  I(\OR, f)/P(\OR, f)
	\end{equation*} 
\end{Prop}

\begin{Rem}
	Si $\OR = \OR_{\K}$, $f = 1$ et aucune restriction n'est faite. 
\end{Rem}

Dans la suite on fixe un ordre $\OR$ de conducteur $f$, et on considère les idéaux de $\OR_{\K}$ de norme premier avec $f$.

\begin{Def}\cite{Cox}[7.C, Lem 7.18]
	Soit $f\in\N$. Un idéal entier $\LR_{\K}\subset\OR_{\K}$ est dit premier avec $f$ si l'une des deux conditions équivalentes suivantes est vérifiées :
	
	\begin{enumerate}
		\item $N(\LR_{\K})$ est premier avec $f$
		\item $\LR_{\K} + f\OR_{\K} = \OR_{\K}$
	\end{enumerate}
	
	On note $I_{\K}(f)$ le sous groupe de $I(\OR_{\K})$ engendré par les idéaux entiers premier avec $f$, et $P_{\K}(f)$ le sous groupe de $I_{\K}(f)$ engendré par les idéaux entiers principaux premier avec $f$
\end{Def}

On montre alors que $I(\OR, f)$ et $I_{\K}(f)$ sont isomorphes.

\begin{The}\cite{Cox}[Prop 7.20]
	
	Soit $\OR$ un ordre d'un corps quadratique imaginaire $\K$, de conducteur $f$.
	\begin{enumerate}
		\item Si $\LR_{\K}\subset\OR_{\K}$ est un idéal premier avec $f$, alors $\LR_{\K}\cap\OR$ est un idéal de $\OR$ premier avec le conducteur, de même norme que $\LR_{\K}$.
		\item Si $\LR\subset\OR$ est un idéal premier avec le conducteur, alors $\LR\OR_{\K}$ est un idéal de $\OR_{\K}$ premier avec $f$, de même norme de que $\LR$.
		\item L'application $\LR\mapsto\LR\OR_{\K}$ induit un isomorphisme entre $I(\OR,f)$ et $I_{\K}(f)$ d'inverse $\LR_{\K}\mapsto\LR_{\K}\cap\OR$
	\end{enumerate}
\end{The}

Une première conséquence de ce théorème est l'application à la factorisation d'idéaux premier avec le conducteur.

\begin{Coro}\cite{Cox}[7.C, Ex 7.26]
	
	Soit $\LR\subset\OR$ premier avec $f$, alors $\OR/\LR \simeq \OR_{\K}/\LR\OR_{\K}$. En particulier l'application $\LR\mapsto\LR_{\K}$ conserve les idéaux premiers. 
\end{Coro}

\begin{Coro}\cite{Cox}[7.C, Ex 7.26]
	Un idéal $\LR\subset\OR$ premier avec le conducteur admet une unique factorisation (à permutation prêt des facteurs) par des idéaux premier, et premiers avec le conducteur.
	
	\begin{equation*}
		\LR = \PR_{1}^{e_{1}}\PR_{2}^{e_{2}}\ldots\PR_{g}^{e_{g}}
	\end{equation*}
	
	avec $e_{1}$, ... $e_{g}\in\N$ et $g\in\N$.
	
\end{Coro}

\begin{Coro}\cite{Cox}[7.C, Ex 7.26]
	
	Tout idéal fractionnaire $\IR \in I(\OR,f)$ admet une unique factorisation par des idéaux entiers premier inclus de $I(\OR,f)$ et par leurs inverses (à permutation prêt des facteurs):
	\begin{equation*}
		\IR = \PR_{1}^{e_{1}}\PR_{2}^{e_{2}}\ldots\PR_{g}^{e_{g}}
	\end{equation*}
	
	avec $e_{1}$, ... $e_{g}\in\Z$ et $g\in\N$.
	
\end{Coro}

Une autre conséquence du théorème précédent est la description du groupe de classe de $\OR$ à partir d'idéaux fractionnaires de $\OR_{\K}$ premier avec $f$.

\begin{Prop}\cite{Cox}[Prop 7.22]
	Si $\OR$ est un ordre d'un corps quadratique imaginaire, de conducteur $f$, alors $Cl(\OR) =  I_{\K}(f)/P_{\K}(f)$
\end{Prop}







\subsection{Ramification des nombres premiers}

Dans $\Z$, l'idéal engendré par un nombre premier $p$ est un idéal premier. Cependant dans un ordre $\OR$ d'un corps de nombre, l'idéal $p\OR$ associé à $p$ peut ne plus être premier. On peut cependant factoriser de façon unique l'idéal $p\OR$ dans deux cas :

\begin{enumerate}
	\item Si $\OR = \OR_{\K}$ est l'ordre maximal d'un corps de nombre $\K$
	\item Si $\OR$ est un ordres d'un corps quadratique imaginaire, et $p$ ne divise pas le conducteur $f$.
\end{enumerate}

On énonce les différentes possibilités de factorisation de $p\OR$:

\begin{Def}
	Soit $\OR$ un ordre d'un corps de nombre $\K$, et $p$ premier, tels que $p\OR$ admette une unique factorisation par des idéaux premier de $\OR$ :
	\begin{equation*}
		p\OR = \PR_{1}^{e_{1}}\PR_{2}^{e_{2}}\ldots\PR_{g}^{e_{g}}
	\end{equation*} 
	
	Si un idéal premier $\PR$ apparaît dans cette factorisation, on dira que $\PR$ est au dessus de $p$. On note $g_{p}=g$ le nombre d'idéaux premiers de la factorisation. Les exposants $e_{i}$ sont appelés les indices de ramification de $p$ dans $\OR$.
	Puisque $p\OR\subset\PR_{i}$ et que $\PR_{i}$ est maximal, $\OR/\PR_{i}$ est un corps de caractéristique $p$ et on note $f_{i}$ son degré. Les coefficients $f_{i}$ sont appelés les degrés d'inerties de $p$ dans $\OR$
\end{Def}

Il est possible d'expliciter les idéaux $\PR_{i}$, d'après un théorème dû à Kummer. On l'énonce d'abord dans le cas des ordre maximaux.

\begin{Prop}\cite{Cox}[Th 5.8, 5.9]
	
	Soit $\K$ un corps de nombre, et $n = \left[ \K : \Q\right] $. Soit $p$ un nombre premier. Alors $N(p\OR) = p^{n}$ et $\sum_{i = 1}^{g}e_{i}f_{i} = n$.
	De plus si l'extension $\K/\Q$ est galoisienne, alors les coefficients $e_{i}$ et $f_{i}$ ne dépendent pas de l'indice $i$ de l'idéal premier $\PR_{i}$ correspondant.  
	On note $e_{p}$ et $f_{p}$ leur valeur. Alors $e_{p}f_{p}g_{p} = n$.
\end{Prop}

\begin{Def}
	Soit $\K$ un corps de nombre qui est une extension galoisienne de $\Q$. On note $n = \left[ \K : \Q\right] $. On dira que $p$ est
	\begin{enumerate}
		\item ramifié dans $\K$, si son indice de ramification $e_{p} > 1$
		\item inerte dans $\K$, si son degré d'inertie $f_{p} = n$, ce qui implique $e_{p} = g_{p} = 1$ et $p\OR_{\K}$ maximal
		\item décomposé dans $\K$, si $g_{p} = n$, ce qui implique $e_{p} = f_{p} = 1$
	\end{enumerate}
\end{Def}

\begin{Rem}
	Si $\K$ est un corps quadratique imaginaire, l'extension $\K/\Q$ est galoisienne, donc pour tout nombre premier $p$, $e_{p}f_{p}g_{p} = 2$ et $p$ est ramifié dans $\K$ si et seulement si $e_{p} = 2$.
\end{Rem}

\begin{The}{Théorème de Kummer}\cite{Coh00}[Th 4.8.13]
	
	
	Soit $\K$ un corps de nombre. Il existe $\alpha\in\OR_{\K}$ tel que $\K = \Q(\alpha)$. Soit $T(X)$ le polynôme minimal de $\alpha$, et $p$ premier ne divisant pas $\left[ \OR_{\K} : \Z\left[ \alpha\right] \right] $. On écrit la factorisation de $T(X)$ sur $\F_{p}$:
	\begin{equation*}
		T(X) = T_{1}(X)T_{2}(X)\ldots T_{r}(X)
	\end{equation*}
	Alors $r = g_{p}$ et $p\OR = \PR_{1}^{e_{1}}\PR_{2}^{e_{2}}\ldots\PR_{r}^{e_{r}}$ avec
	\begin{equation*}
		\PR_{i} = p\OR + T_{i}(\sqrt{\alpha})\OR 
	\end{equation*}
	De plus $deg(T_{i}) = f_{i}$
\end{The}

\begin{Rem}
	Le théorème 1.20 et le corollaire 1.22 permettent de généraliser ce résultat à tout ordre $\OR$ d'un corps quadratique imaginaire $\K = \Q(\alpha)$ tel que $\Z\left[ \alpha\right] \subset \OR$, car dans ce cas le conducteur de $\OR$ divise $\left[ \OR_{\K} : \Z\left[ \alpha\right] \right]$ et $p\OR$ sera premier avec le conducteur. 
\end{Rem}

On énonce donc le théorème de Kummer dans le cas des corps quadratique imaginaire. 

\begin{Prop}{Théorème de Kummer, cas quadratique imaginaire}
	
	Soit $\K$ un corps quadratique imaginaire, et $\OR$ un ordre de $\K$ de discriminant $\Delta$. Alors $\K = \Q(\sqrt{\Delta})$. Soit $T(X) = X^2 - \Delta \in\F_{p}\left[ X\right] $. Alors
	\begin{enumerate}
		\item Si $T(X) = X^2$, c'est à dire $p$ divise $\Delta$, alors $p$ est ramifié.
		\begin{equation*}
			p\OR = \PR^{2}\; avec \; \PR = p\OR + \sqrt{\Delta}\OR
		\end{equation*}
		
		\item Si $T(X)$ est irréductible, alors $p$ est inerte, $p\OR$ est un idéal premier de norme $p^{2}$.
		
		\item Si $T(X) = T_{1}(X)T_{2}(X)$ se factorise de façon non trivial, alors $p$ est décomposé.
		\begin{equation*}
			p\OR = \PR_{1}\PR_{2}\; avec \; \PR_{i} = p\OR + T_{i}(\sqrt{\Delta})\OR
		\end{equation*}
		
	\end{enumerate}
\end{Prop}

Ainsi connaître le comportement de $p$ dans un ordre quadratique imaginaire $\OR$ de discriminant $\Delta$ se résume à calculer le symbole de Legendre $ \left( \frac{\Delta}{p}\right) $. Expliciter le cas décomposer demande en plus de trouver les racines de $\Delta$ dans $\F_{p}$. 

\subsection{Corps de classe de Hilbert}

Soit $\K$ un corps de nombre, $\Lc/\K$ une extension finie. On s'est intéressé au comportement des entiers premiers $p$ dans un l'ordre maximal $\OR_{\K}$. On peut de même s'intéresser aux comportement des idéaux premier de $\OR_{\K}$ vu dans $\OR_{\Lc}$. On commence par généraliser aux idéaux le vocabulaire introduit sur les nombres premiers.

\begin{Def}
	Soit $\K$ un corps de nombre, et $\PR_{\K}\subset\OR_{\K}$ un idéal premier. Soit $\Lc/\K$ une extension galoisienne. L'idéal $\PR_{\K}\OR_{\Lc}$ se factorise dans $\OR_{\Lc}$ :
	\begin{equation*}
		\PR_{\K}\OR_{\Lc} = \PR_{1}^{e_{1}}\PR_{2}^{e_{2}}\ldots\PR_{g}^{e_{g}}
	\end{equation*} 
	On note $g$ le nombre d'idéaux premiers de la factorisation. Les exposants $e_{i}$ sont appelés les indices de ramification de $\PR_{\K}$ dans $\Lc$.
	Puisque $\PR_{\K}\OR_{\Lc}\subset\PR_{i}$ et que $\PR_{i}$ est maximal, le corps $\OR_{\Lc}/\PR_{i}$ est une extension du corps $\OR_{\K}/\PR_{\K}$ et on note $f_{i}$ son degré. Les coefficients $f_{i}$ sont appelés les degrés d'inerties de $\PR_{\K}$ dans $\Lc$
\end{Def}


\begin{Prop}\cite{Cox}[Th 5.8, 5.9]
	
	Soit $\K$ un corps de nombre, $\Lc$ une extension galoisienne de $\K$ et $n = \left[ \Lc : \K\right] $. Soit $\PR_{\K}\subset\OR_{\K}$ un idéal premier. Alors les coefficients $e_{i}$ et $f_{i}$ de $\PR_{\K}$ dans $\Lc$ ne dépendent pas de l'indice $i$ i.e. de l'idéal premier $\PR_{i}\subset\OR_{\Lc}$ correspondant.  
	On note $e$ et $f$ leur valeur. Alors $efg = n$.
\end{Prop}

\begin{Def}
	Soit $\K$ un corps de nombre, $\Lc$ une extension galoisienne de $\K$. On note $n = \left[ \Lc : \K\right] $. Un idéal premier $\PR_{\K}\subset\OR_{\K}$ est dit :
	\begin{enumerate}
		\item ramifié dans $\Lc$, si son indice de ramification $e > 1$
		\item inerte dans $\Lc$, si son degré d'inertie $f = n$, ce qui implique $e = g = 1$ et $p\OR_{\K}$ maximal
		\item décomposé dans $\Lc$, si $g = n$, ce qui implique $e = f = 1$
	\end{enumerate}
\end{Def}

Le but de ce qui suit est d'introduire le une extension particulière de $\K$, appelée son corps de classe de Hilbert, dont les propriétés sont liées au groupe de classe de $\OR_{\K}$.

\begin{Def} Soit $\Lc/\K$ une extension galoisienne d'un corps de nombre  $\K$. Une injection $\sigma :\K\rightarrow\C$ est dite réelle si elle plonge $\K$ dans $\R$. Une injection réelle $\sigma$ de $\K$ est dite ramifiée dans $\Lc$ si elle peut se prolonger en une injection $\sigma :\Lc\rightarrow\C$ qui n'est pas réelle. 
\end{Def}

\begin{Rem}
	Les injections $\sigma :\K\rightarrow\C$ sont appelées les "premiers infinis" de $\K$. \cite{Cox}[5.C]
\end{Rem}

\begin{Def}
	Une extension $\Lc/\K$ galoisienne d'un corps de nombre $\K$ est dit abélienne si son groupe de Galois est abéliens. Une tel extension est non-ramifié si ni aucun idéal premier de $\OR_{\K}$ ni aucune injection réelle de $\K$ n'est ramifié dans $\Lc$.
\end{Def}

\begin{The}\cite{Cox}[Th 5.18]
	
	Soit $\K$ un corps de nombre. Il existe une unique extension $\HR/\K$ abélienne, non-ramifiée et maximale au sens de l'inclusion, c'est à dire telle que toute extension abélienne non-ramifiée de $\K$ est incluse dans $\HR$
\end{The}

Cette extension de $\K$ est liée au groupe de classe de $\OR_{\K}$ par le théorème de réciprocité d'Artin. L'énoncé exacte de ce théorème demande la construction du symbole d'Artin, qui généralise le symbole de Legendre pour les idéaux premiers \cite{Cox}[Lem 5.19, Cor 5.21]. 

\begin{The}{Théorème de réciprocité d'Artin pour le corps de classe}\cite{Cox}[Th 5.23]
	
	Soit $\K$ un corps de nombre, $\HR$ son extension abélienne non-ramifiée maximale. Alors le groupe de Galois de l'extension $\HR/\K$ est isomorphe au groupe de classe $Cl(\OR_{\K})$
	\begin{equation*}
		Gal(\HR/\K)\simeq Cl(\OR_{\K})
	\end{equation*}
\end{The} 

\begin{Def}
	Soit $\K$ un corps de nombre. Son extension abélienne non-ramifiée maximale $\HR$ est appelé corps de classe de Hilbert de $\K$.
\end{Def}